\section{An explicit sign unit inside the smaller cyclic host}
\label{sec:cyclic-cm-unit}

The larger-field unit question isolated in AC-G32 has a
constructive answer.  Two conjugates of one cyclotomic norm
suffice.  The norm is taken before the unit is used, so the
field of the smaller Host is unchanged.

\begin{theorem}[AC-G33: the missing signature unit and elliptic principal model]
\label{AC-G33}
Let $d$ be as in AC-G32 and put $D=3d$ and
$R=\mathbb Q(\zeta_{16})^+$.  There is an explicitly
specified integral unit
\[
 \kappa_d\in\mathcal O_{R(\sqrt D)}^\times
              \subset\mathcal O_{E_h^+}^\times
\]
with
\[
 \operatorname{sign}\tau(\kappa_d)
     =\operatorname{sign}\tau(\sqrt D)
                 \quad\text{at every real embedding},        \tag{G33.1}
\]
and absolute logarithmic Weil height
\[
             h(\kappa_d)\leq2(d-1)\log2+\log105.               \tag{G33.2}
\]
It requires only a cyclotomic norm, one fixed conjugate,
and at most one copy of each of three fixed circular units,
with an optional minus sign.

Consequently $\beta=\xi_0/\kappa_d$ makes the unchanged
source lattice $AA$ of AC-G32 a principally polarized
CM lattice of type induced from $K=\mathbb Q(\omega)$.
With $g=[E_h:K]=2^{h+2}$, its unpolarized abelian variety is
\[
                       C_0^g,\qquad C_0:y^2=x^3+1.             \tag{G33.3}
\]
The induced integral $\mathcal O_{E_h}$-action and the
principal polarization can be realized on this power over
$K$.  The polarization need not be its product polarization.

For any fixed $A,a\geq0$, the source window bound
$\log d\leq Y_0^{\epsilon/5}(\log Y_0)^2$, with
$x\geq Y_0\geq e$ and fixed $\epsilon>0$, gives
\[
        x^A(1+h(\kappa_d))^a
                    =o\bigl(\exp(x^\epsilon/2)\bigr).
                                                               \tag{G33.4}
\]
This controls the displayed logarithmic height and fixed
polynomial expressions in it.  It does not identify that
quantity with an unlogged theta norm or prove the mixed-block
arithmetic tie.
\end{theorem}
\begin{proof}
\proofstep{1. A real unit from the quadratic cyclotomic norm.}
Put $t=\zeta_{16}$ and $z=\zeta_D$.
The primitive quadratic character
\[
        \chi_D(a)=(a/3)(a/d)
                 \quad(a\in(\mathbb Z/D\mathbb Z)^\times)
\]
is even and cuts out $\mathbb Q(\sqrt D)$.
Indeed it is the product of the characters of quadratic
discriminants $-3$ and $-d$, whose compositum has this real
quadratic subfield.  The fields of conductors $16$ and $D$
are linearly disjoint.  Define
\[
 \begin{split}
 P&=\prod_{\chi_D(a)=1}(1-tz^a)
    =N_{\mathbb Q(t,z)/\mathbb Q(t,\sqrt D)}(1-tz),\\
 v&=t^{-(d-1)/2}P.
 \end{split}\tag{G33.5}
\]
The product has $d-1$ factors.  Each $tz^a$ is a primitive
$16D$-th root, and $1-tz^a$ is an integral unit because
$\Phi_{16D}(1)=1$.
For this last equality one may use
$\Phi_{16D}(X)=\Phi_{2D}(X^8)$ and
$\Phi_{2D}(1)=\Phi_D(-1)=\Phi_d((-1)^3)/\Phi_d(-1)=1$.
Thus both $P$ and $v$ are units.

The sum of the residues $a$ occurring in (G33.5) is zero
modulo $D$.  Modulo $d$ they run once through every nonzero
residue: its quadratic sign determines the unique nonzero
residue modulo three.  Modulo three there are $(d-1)/2$
of each of its two nonzero residues, since the two quadratic
signs modulo $d$ occur equally often.  Both sums are zero.
As $d-1$ is even, complex conjugation consequently gives
\[
 \overline P=(-t^{-1})^{d-1}P=t^{-(d-1)}P,
                       \qquad\overline v=v.          \tag{G33.6}
\]
The fixed real field of $\mathbb Q(t,\sqrt D)$ is
$R(\sqrt D)$, so $v$ lies in precisely this field.
The auxiliary root $z$ has disappeared through the norm;
no extension of the Host has been made.

\proofstep{2. Two conjugates give opposite signs in every quadratic pair.}
Let $\sigma_D$ fix $t$ and negate $\sqrt D$.
It exchanges the two character cosets in (G33.5), hence
\[
 n:=v\sigma_D(v)=t^{-(d-1)}\Phi_D(t).                 \tag{G33.7}
\]
The polynomial identity
\[
 \Phi_{3d}(X)
       =\frac{(X^{3d}-1)(X-1)}{(X^3-1)(X^d-1)}
\]
and $d\equiv7\pmod {16}$ show that, at the real embedding
of $R$ indexed by $b\in\{1,3,5,7\}$,
\[
 n_b=-\frac{\sin(5b\pi/16)\sin(b\pi/16)}
                  {\sin(3b\pi/16)\sin(7b\pi/16)}.
                                                               \tag{G33.8}
\]
For example, substituting $t^d=t^7$ and $t^{3d}=t^5$
in (G33.7) leaves the phase $t^{-8}=-1$;
this fixes its sign without a choice of square root.
None of the displayed sines vanishes, and their signs give
\[
 \begin{array}{c|rrrr}
 b&1&3&5&7\\\hline
 \operatorname{sign}(n_b)&-&+&+&-
 \end{array}.                                                   \tag{G33.9}
\]
Let $\tau_3$ send $t$ to $t^3$ and fix $\sqrt D$.
On the four real embeddings it permutes the indices by
multiplication by three modulo $\{\pm1\}$.
Inspection of (G33.9) shows
$n_b n_{3b}<0$ at all four embeddings.
Therefore the unit
\[
                             w=v\tau_3(v)                      \tag{G33.10}
\]
has totally negative norm from $R(\sqrt D)$ to $R$.
Its two real values above every embedding of $R$ have
opposite signs.

\proofstep{3. Aligning those four pairs by fixed circular units.}
For $a=3,5,7$, put
\[
 c_a=t^{(1-a)/2}\frac{1-t^a}{1-t}\in\mathcal O_R^\times.
                                                               \tag{G33.11}
\]
The quotient is an integral geometric sum.  Its inverse
is integral as well: choose $b$ with $ab\equiv1\pmod {16}$
and use $(1-(t^a)^b)/(1-t^a)$, adjusting by a root of unity.
Conjugation fixes the expression, so it is a real unit.
Its values at the four real embeddings are the sine ratios
$\sin(ab\pi/16)/\sin(b\pi/16)$, with signs
\[
 \begin{array}{c|rrrr}
 &1&3&5&7\\\hline
 -1&-&-&-&-\\
 c_3&+&+&+&-\\
 c_5&+&+&-&+\\
 c_7&+&-&+&-
 \end{array}.                                                   \tag{G33.12}
\]
These four sign vectors form a basis over $\mathbb F_2$.
This is the same explicit four-dimensional calculation as
AM-H16, (H16.7)--(H16.9); no general unit-signature theorem
is needed.

Above each $b$, take the extension sending $\sqrt D$ to
its positive real value, and let $a_b\in\{0,1\}$ record
whether $w$ is negative there.  Reduce the following
exponents modulo two and set
\[
 w_0=(-1)^{a_1}c_7^{a_1+a_3}c_5^{a_1+a_5}c_3^{a_3+a_7},
                  \qquad \kappa_d=w/w_0.             \tag{G33.13}
\]
The table makes $w_0$ and $w$ have the same sign at
these four selected extensions.  On the other member of
each pair $w_0$ is unchanged and $w$ reverses sign by
(G33.10).  This proves (G33.1).
All sign tests involve nonzero real algebraic numbers,
so exact real-root isolation specifies this construction
without a numerical choice.

There are $d-1$ factors of height at most $\log2$ in
(G33.5); roots of unity have height zero.  Thus
$h(v)\leq(d-1)\log2$ and
$h(w)\leq2(d-1)\log2$.
Each $c_a$ is a sum of $a$ roots of unity, giving
$h(c_a)\leq\log a$.  Formula (G33.13) uses at most one
of each, so $h(w_0)\leq\log105$.
Subadditivity of height under multiplication and
$h(w_0^{-1})=h(w_0)$ prove (G33.2).

\proofstep{4. The unchanged principal lattice and the elliptic power.}
The field $R(\sqrt D)$ is contained in (G32.9).
Apply AC-G32-signature to the constructed unit.
For every embedding $\phi$ in the induced $K$-type,
write $\phi(\beta)=ib_\phi$ with $b_\phi>0$.
The alternating form
\[
                     (x,y)\longmapsto
                  \operatorname{Tr}_{E_h/\mathbb Q}(\beta x\bar y)
\]
is integral and unimodular by (G32.5), and its value on
$(x,ix)$ is $2\sum_\phi b_\phi|x_\phi|^2>0$.
These are the Riemann relations, giving the principal
polarization on $\mathbb C^{\Phi}/\Phi(AA)$.
Multiplication by $\mathcal O_{E_h}$ preserves $AA$;
the trace identity shows that its Rosati involution is
complex conjugation.  This is exactly the CM lattice
argument of AM-H17, step~6, with its primary locators
Milne, \emph{Complex Multiplication}, Example~2.9,
pp.~25--26, and Proposition~3.17, p.~31.

The Eisenstein ring $\mathcal O_K=\mathbb Z[\omega]$
is norm-Euclidean: rounding first the imaginary coordinate
and then the real coordinate leaves squared absolute value
at most $3/16+1/4<1$.  Consequently the torsion-free
finite module $AA$ is free over $\mathcal O_K$, of rank
$g=[E_h:K]$.
An $\mathcal O_K$-basis identifies the complex lattice of
its induced type with $\mathcal O_K^g$, up to the common
complex scaling that identifies $\mathbb C/\mathcal O_K$
with $C_0$.  This proves (G33.3).
After transporting the polarization and the action to
$C_0^g$, their homomorphisms have entries in
$\operatorname{End}_{\mathbb C}(C_0)=\mathcal O_K$.
Every such endomorphism is defined over $K$, generated by
integer multiplications and $(x,y)\mapsto(\omega x,y)$.
Thus the transported action and polarization morphism are
defined over $K$.  Positivity and principality can be checked
after embedding in $\mathbb C$, so this is the asserted
principal abelian model.  This reasoning does not identify
its polarization with the product one.

\proofstep{5. The quantitative statement has logarithmic normalization.}
For $x\geq Y_0\geq e$, the given upper bound on $\log d$
is at most $x^{\epsilon/5}(\log x)^2$.
For fixed $A,a$, (G33.2) therefore gives
\[
 \log\bigl(x^A(1+h(\kappa_d))^a\bigr)
     \leq A\log x+a x^{\epsilon/5}(\log x)^2+O_a(1)
                         =o(x^\epsilon).
\]
Subtracting $x^\epsilon/2$ proves (G33.4).
The assertion is about this explicitly normalized height,
not its exponential or a separately defined theta quantity.
\end{proof}

\paragraph{Source scope after this construction.}
The signature-unit existence gate for AC-G32 is now closed.
The proof gives a principal elliptic CM model on the exact
source lattice, with a controlled logarithmic sign-unit height,
inside a host that is one literal \texttt{IsChiField}.
It still changes the CM type from the source type induced
from $K_d$ to the one induced from $K$.
The source ordinary clause remains a statement about the
former type.  For instance, at eleven the latter elliptic
curve has trace zero: cubing permutes $\mathbb F_{11}$,
so its quadratic-character point-count sum is zero.
The ordinary and theta identities for the new type therefore
require their actual local hypotheses and comparison proofs.
The cross-block character tie, the main cubic-twisted
cohomology supply, and its exact theta normalization are
not consequences of the sign unit or of (G33.4).

\paragraph{Algebraic-geometric dictionary.}
The principal lattice statement uses the audited CM lattice
and Riemann-form dictionary in Milne,
\emph{Complex Multiplication}, Example~2.9, pp.~25--26,
and Proposition~3.17, p.~31, also used in AM-H17.
The finite character field remains the finite etale-cover
object of SGA~1, Expos\'e~V, \S\S4 and 7, after inversion
of its conductor primes.  The explicit unit supplies the
additional integral polarization data; it does not identify
a finite-cover coefficient with a theta or cohomology
comparison.  These are the already audited locators;
no new Grothendieck archive cote is asserted.

\begin{corollary}[AC-G33-twist: two principal scalar cubic packets]
\label{AC-G33-twist}
The principal model $A_0=C_0^g$ of AC-G33 has a scalar
$\mu_3$-action preserving its polarization and commuting
with its integral $\mathcal O_{E_h}$-action.
Put $p=269$, $C_p:y^2=x^3+p^2$ and $A_p=C_p^g$.
There is a principal polarization on $A_p$, defined over
$K$, and the same integral $\mathcal O_{E_h}$-action,
with the following properties.  If $\alpha^3=p$ and
$\chi(\sigma)=\sigma(\alpha)/\alpha$, then the product
isomorphism
\[
 \phi_g:A_0\longrightarrow A_p,
       \qquad (x,y)\longmapsto(\alpha^2x,py)
                  \quad\text{on each factor}                 \tag{G33.14}
\]
is $\mathcal O_{E_h}$-linear and polarized over $K(\alpha)$.
Its Galois conjugates satisfy
\[
                  \sigma(\phi_g)=\phi_g\,[\chi(\sigma)^2].
                                                               \tag{G33.15}
\]
Consequently every covariant Tate module is twisted by
$\chi^2$ integrally, including at coefficient prime three.
Its reverse $\mathcal O_{E_h}$-linear Hom is a rank-one
integral coefficient with character $\chi$.
The two principal varieties have identical geometric CM
data and stable Faltings heights.  This construction invokes
no Fourier projector and no full regular tensor.
\end{corollary}
\begin{proof}
Identify $A_0$ with $C_0^g$ as in AC-G33, and use the
product principal polarization only to identify its dual.
The constructed principal polarization is represented by
an integral positive Hermitian matrix
$H\in M_g(\mathcal O_K)$.
Every scalar $u\in\mu_3$ satisfies $\bar u u=1$, so
$\bar u H u=H$.  This proves the asserted polarization
invariance; the scalar action commutes with all matrices
representing $\mathcal O_{E_h}$.

On $C_p$, the same Eisenstein endomorphisms are defined over
$K$, and (G33.14) commutes with them.  An elliptic group
isomorphism preserves the principal polarization defined by
the origin.  Therefore the same matrix $H$ on $C_p^g$
pulls back under $\phi_g$ to the polarization on $C_0^g$.
It is defined over $K$; positivity and degree one are
checked after this isomorphism.  The same integral action
matrices also commute with $\phi_g$, proving the statement
about $\mathcal O_{E_h}$.  Taking Galois conjugates of the
coordinate formula gives (G33.15) directly.

For a covariant Tate vector $v$, formula (G33.15) reads
$\sigma(\phi_g(v))=\phi_g(\chi(\sigma)^2\sigma v)$.
This proves the integral twist, without division by three.
The original Tate lattice is rank one projective over
$\mathcal O_{E_h}\otimes\mathbb Z_\ell$: the characteristic-zero
integral comparison identifies it with the invertible
$\mathcal O_{E_h}$-ideal $AA$ tensored with $\mathbb Z_\ell$.
The latter is locally free of rank one over this product
of local rings.  Thus its endomorphism coefficient is
$\mathcal O_{E_h}\otimes\mathbb Z_\ell$; reversing the Hom
replaces $\chi^2$ by $\chi$.  This also fixes the variance
at $\ell=3$.
Finally the varieties are geometrically isomorphic, which
preserves their CM data and stable Faltings heights.
This is the concrete scalar instance of AC-G24 on the
principal model just constructed.
\end{proof}

The corollary supplies separate principal scalar packets.
A common twist within the one distinguished source Speh
block still cancels from its pairwise Hom coefficients,
as proved in AC-G29.  To use the nontrivial Hom between
these two packets in the source bound requires the actual
cross-block theta and cohomology comparison.  The polarized
isomorphism above does not assert that comparison or a
comparison of ramified integral theta lattices.
